% Solutions to the Chapter 12 exercises that are not code, from lectures/L12/appendix/c_exercises.md.
% Exercises 12.6 and 12.7 are Code and are not answered; see chapter.tex for why.

\section{Chapter 12: Real-Time Linux}
\label{sol:sec:c12}
Answers to the exercises in \secref{c12:sec:exercises}. Exercises~\ref{c12:ex:6} and~\ref{c12:ex:7}
are Code, and are checked by running them on the target.

% ------------------------------------------------------------------------------------------
\solution[fillaccent2]{c12:ex:1}{What real time claims}{Recall}

\xpart{a} A real-time system is one whose answer is correct only if it arrives within a stated
deadline, every time: a late answer is a wrong answer, however right its content.

\xpart{b} \textbf{For 6 ms, B is real time and A is not.} B's worst observed response, 5.1 ms, is
inside the deadline; A's is 30 ms, five times the deadline, and its 50 us average does not change
that. \textbf{For 1 ms, neither is.} B misses every time, since it never responds in less than
5 ms. A meets the deadline on most responses but not on all of them, and a deadline met most of the
time is a deadline missed. Being real time is a property of a system against a deadline, not of the
system alone.

\xpart{c} ``Which number improved: the average or the worst case?'' A mean 30\% lower says nothing
about the maximum, and a real-time requirement is a statement about the maximum
(\secref{c12:app:a1}). The next questions are under what load and for how long
(\secref{c12:app:b5}).

\xpart{d} In \textbf{hard} real time a missed deadline is a system failure: an airbag, a motor
commutation loop. In \textbf{soft} real time it degrades quality: audio, video. Linux with
\code{PREEMPT_RT} is honestly good soft real time. It is used for hard real time, but only by
people who have measured their own system very carefully, because nobody can prove a worst case
for a kernel of that size; it can only be measured. It is not a hard real-time operating system in
the sense that a small RTOS with a provable worst case is.

% ------------------------------------------------------------------------------------------
\solution[fillaccent2]{c12:ex:2}{The three sources}{Recall}

\xpart{a} Interrupts disabled, preemption disabled, and priority inversion (\secref{c12:app:a2}).

\xpart{b}
\begin{itemize}
  \item \textbf{A \code{spin_lock_irqsave} section copying 4 KB:} interrupts disabled, on that CPU,
    for the length of the copy, with preemption disabled along with them. (On \code{PREEMPT_RT} a
    \code{spinlock_t} taken with \code{_irqsave} does not disable interrupts at all,
    \file{Documentation/locking/locktypes.rst}:238. It is a sleeping lock there, and the cost of
    the section moves to whoever waits for the lock.)
  \item \textbf{A \code{udelay(200)} in a probe function:} on this course's \code{PREEMPT} kernel,
    none of the three by itself. \code{probe} runs in process context and is preemptible, so a
    higher-priority task simply preempts it, and the 200 us is lost only to the probe. It becomes
    preemption-disabled time on a \code{PREEMPT_NONE} or \code{PREEMPT_VOLUNTARY} kernel, where
    there is no preemption point inside it, or on any kernel if a spinlock is held around it, and
    interrupts-disabled time under \code{spin_lock_irqsave}. That is why busy-waiting is a habit of
    its own in \secref{c12:app:a6}: what it costs depends on the context around it.
  \item \textbf{A low-priority task holding a mutex a high-priority task needs:} priority inversion,
    as soon as a medium-priority task keeps the holder off the CPU. Without the medium task it is
    ordinary blocking, bounded by the holder's critical section.
  \item \textbf{A hard interrupt handler that takes 300 us:} interrupts disabled. A hard handler runs
    with interrupts masked on its CPU, so for 300 us no other interrupt is taken there and nothing
    is scheduled there either.
\end{itemize}

\xpart{c} \textbf{Interrupts disabled and preemption disabled are somebody's code}: a section someone
wrote, fixed by making it shorter. \textbf{Priority inversion is structural}: it needs no long
section at all, only three tasks, one lock and an unlucky interleaving. The mechanism that addresses
it is \textbf{priority inheritance}: while a low-priority task holds a lock a higher-priority task
waits for, it runs at the waiter's priority, so no medium-priority task can preempt it
(\secref{c12:app:a3}). In Linux that is \code{rt_mutex}, which \code{PREEMPT_RT} puts under
\code{mutex} and \code{spinlock_t}.

\xpart{d} Mars Pathfinder landed in July 1997, and a few days later began suffering total system
resets, each of which lost data. The high-priority bus management task needed a mutex held by the
low-priority meteorological task, and the medium-priority communications task, which needed
neither, preempted the meteorological task and kept it from releasing the mutex. The bus management
task therefore did not run in time, and a watchdog concluded that something had gone badly wrong
and reset the computer. The VxWorks mutex had been created with priority inheritance switched off,
and the fix, uploaded from Earth, switched it on.

% ------------------------------------------------------------------------------------------
\solution[fillaccent3]{c12:ex:3}{A latency budget}{Hand calculation}

\xpart{a} Add the stages:
\[
  40 + 15 + 120 + 60 + 200 = 435~\text{us} \le 500~\text{us}.
\]
It meets the deadline, with 65 us to spare. The sum of the per-stage maxima is a safe bound on the
chain only because each figure is a maximum; the stages' worst cases need not happen together, so
the true end-to-end worst case is at most 435 us.

\xpart{b} With the work in the hard handler, the reaction happens there, and nothing on the path
waits for the thread or for userspace. The handler now does its own 15 us and the thread's 60 us:
\[
  40 + (15 + 60) = 115~\text{us}.
\]
It meets the deadline with 385 us to spare. What it has done to every other device: the hard
handler now runs for 75 us with interrupts masked on its CPU, so every other interrupt routed there
can wait up to 75 us longer, and no task of any priority runs there meanwhile. The deadline was met
by moving 75 us into everybody else's worst case, the first of the three sources in
\secref{c12:app:a2}. (On \code{PREEMPT_RT} the kernel would force this handler into a thread unless
it asked for \code{IRQF_NO_THREAD}, which it would then need, and which makes the imposition
explicit.)

\xpart{c} The other driver's \code{spin_lock} disables preemption, not interrupts. It therefore
delays the points at which a task is scheduled, the threaded handler and the reader, and not the
hard interrupt.
\begin{itemize}
  \item \textbf{a)} $435 + 300 = 735$ us. It fails, by 235 us.
  \item \textbf{b)} 115 us, unchanged: a preemption-disabled section cannot delay a hard interrupt,
    because interrupts are still enabled. It survives.
\end{itemize}
\textbf{The arrangement in b) survives.} Even if the other driver had taken its lock with
\code{spin_lock_irqsave}, b) would be $115 + 300 = 415$ us and would still survive, while a) stays
at 735 us.

\xpart{d} Back to $40 + 15 + 120 + 60 + 200 = 435$ us, meeting the deadline with 65 us to spare.
\textbf{What happened:} the other driver's \code{spinlock_t} is an \code{rt_mutex} under
\code{PREEMPT_RT}, so its holder is preemptible. Your threaded handler runs at \code{SCHED_FIFO}
priority 50 by default (\code{sched_set_fifo}, \file{kernel/sched/syscalls.c}:858), preempts a
holder that is an ordinary task, and your chain stops paying for somebody else's critical section.
\textbf{What it has cost you:}
\begin{itemize}
  \item Every stage figure was measured on \code{PREEMPT} and no longer holds. Under
    \code{PREEMPT_RT} the averages rise: in this chapter's measurement the mean rose by 23\% idle
    and 42\% loaded (\secref{c12:app:a8}). Your own hard handler, requested without
    \code{IRQF_NO_THREAD}, is itself forced into a thread, so the 15 us and 120 us stages change
    shape. The 435 us means nothing until it has been measured again on the real-time kernel.
  \item The other driver's section now takes longer in wall-clock time, because it is preempted.
  \item If your path needs the same lock, you still wait for its holder, boosted by inheritance,
    for up to its 300 us. If the holder is another interrupt thread at the same priority 50,
    \code{SCHED_FIFO} does not preempt an equal, and the 300 us stays in your budget until you raise
    your thread's priority. Priorities are now something you have to design.
\end{itemize}

\xpart{e} \textbf{Which numbers:} no longer the maximum of each stage but the distribution of the
end-to-end latency, and in particular its 99.9th percentile, measured end to end. Percentiles do
not add: the 99.9th percentile of a sum is not the sum of the stages' 99.9th percentiles, and can be
larger or smaller, so a budget of per-stage figures no longer works. The requirement must also say
what the other 0.1\% may be. At one interrupt a millisecond, one in a thousand is one late reaction
every second, so it needs to say how late one may be and whether two may come in a row.

\textbf{Easier to demonstrate}, because a percentile can be established by a finite measurement
with a stated confidence and a maximum never can. If $n$ independent samples show no miss, the
miss rate is below 0.1\% with 95\% confidence once $(1 - 0.001)^n \le 0.05$, which is
$n \ge \ln 0.05 / \ln 0.999 \approx 2{,}995$: about three thousand samples, and more for every miss
you do see. Latency spikes are not independent, since one blocking event delays a run of consecutive
interrupts, so in practice you need many times that, under the product's load. It is still a
question a measurement can answer. Whether the requirement is maximum or percentile, nothing can
be demonstrated from per-stage figures alone.

% ------------------------------------------------------------------------------------------
\solution[fillaccent4]{c12:ex:4}{Reading your own driver for hostility}{Design}

The answer below is for a driver written to the specifications of Exercises~\ref{c8:ex:6},
\ref{c9:ex:6}, \ref{c10:ex:6} and~\ref{c11:ex:6}. Its handler is registered with
\code{devm_request_irq} and no flags. It reads \code{IRQ_STATUS}, drains the device's FIFO, keeps
the latest sample in an \code{atomic_t} for \code{read_raw}, and acknowledges. It also has a
\code{gpio_chip} over \code{GPIO_OUT} and \code{GPIO_IN}. Yours may differ, and the method is the
point.

\xpart{a} \textbf{The hard handler, for its whole length.} It runs with interrupts masked on its
CPU while it reads \code{IRQ_STATUS}, drains up to \code{QA_DEV_FIFO_DEPTH}, 16 samples, one MMIO
read each, and writes the acknowledgement: bounded by about twenty register accesses, each a trap
into QEMU's device model on this target. \textbf{Any \code{spin_lock_irqsave} in process context.}
Chapter~\ref{c9:ch}'s \code{read} took the FIFO lock that way (\secref{c8:app:a8}).
Chapter~\ref{c11:ch} deleted \code{read}, so with \code{read_raw} using \code{atomic_read} as \secref{c11:app:a4}'s does,
none remains, unless your gpiochip's \code{set} does its read-modify-write of \code{GPIO_OUT} under
\code{spin_lock_irqsave}, which is a few register accesses long.
\emph{Under \code{PREEMPT_RT}:} the handler is forced into a thread, because it was requested with
neither \code{IRQF_NO_THREAD} nor \code{IRQF_ONESHOT} (\file{kernel/irq/manage.c}:1372). It runs
as \code{irq/N-...} with interrupts enabled, and the interrupts-off time shrinks to the kernel's
generic primary handler. \code{spin_lock_irqsave} on a \code{spinlock_t} no longer disables
interrupts at all. The cost is nothing and the fix is to leave it alone. What must not be added is
\code{IRQF_NO_THREAD} to ``keep it fast''.

\xpart{b} \textbf{The FIFO lock in the handler}, if it is held across the drain loop: up to 16 MMIO
reads under a lock. Chapter~\ref{c9:ch}'s \code{read} also copied out of the driver's FIFO under
it; that went with \code{read} in Chapter~\ref{c11:ch}, unless you kept the FIFO. The fix is to read
the samples into a local array first and take the lock only to copy them in.
\emph{Under \code{PREEMPT_RT}:} the lock is an \code{rt_mutex}, so the section stops being
preemption-disabled time. Instead it is time for which a waiter of any priority can be blocked,
bounded by inheritance, and the handler thread and a process-context holder now sleep on each other
rather than spin. Shorten the section; do not change its type. A \code{raw_spinlock_t} here would
put back exactly the latency \code{PREEMPT_RT} removed (\secref{c12:app:a5}).

\xpart{c} \textbf{None.} Nothing in Chapters~\ref{c8:ch} to~\ref{c11:ch} asks for \code{udelay},
\code{mdelay}, \code{ndelay} or a loop that polls a register until it changes. The drain loop is
not a busy-wait: it reads until the FIFO is empty, and the device's depth of 16 bounds it.
\emph{Established by} searching the source for \code{udelay}, \code{mdelay}, \code{ndelay},
\code{cpu_relax} and the \code{readl_poll_timeout} family, and by reading every loop whose exit
condition is a register value to check that it has a bound other than the device's cooperation. Had
there been one, the fix is \code{usleep_range} in process context (\secref{c9:app:b3}). Under
\code{PREEMPT_RT} a \code{udelay} in the threaded handler burns its time preemptibly, but it still
holds up the thread and everything waiting for it.

\xpart{d} \textbf{None.} The counters are \code{atomic_t} (Exercise~\ref{c8:ex:6}), the sample is an
\code{atomic_t}, and there is no \code{DEFINE_PER_CPU}, \code{this_cpu_}, \code{smp_processor_id},
\code{get_cpu} or \code{preempt_disable}. \emph{Established by} searching for exactly those. The
nearest thing is the handler's plain \code{spin_lock}, which \secref{c8:app:a8} justifies by
``interrupts are already masked there''. Under \code{PREEMPT_RT} that reason no longer holds, since
the handler is a thread with interrupts enabled, but the code is still correct. The process-context
side takes the same \code{spinlock_t}, both are now the same sleeping lock, and forced threading
adds \code{IRQF_ONESHOT} (\file{kernel/irq/manage.c}:1382), so the line stays masked until the
thread finishes and the handler cannot run twice at once.

\xpart{e} \textbf{Long critical sections}, and interrupts-off time generally. Reading shows where a
section starts and ends, but not how long it takes: a loop's bound comes from the device at run
time, and a call into unfamiliar code has a cost the call site does not show. Use the kernel's
latency tracers, which \file{kernel/qa.config} enables for this chapter and both kernels have.
\code{irqsoff} and \code{preemptoff} (\code{CONFIG_IRQSOFF_TRACER}, \code{CONFIG_PREEMPT_TRACER})
record the longest interrupts-off or preemption-off section since they were reset, with the stack
that caused it. \code{wakeup_rt} (\code{CONFIG_SCHED_TRACER}) records the longest wait of the
highest-priority task to get the CPU. An answer that picks per-CPU assumptions has a case too: the
reliance is an absence, and nothing in the code names it. It must then say what finds it, and no
tool does. You find it by reading every per-CPU access and asking what keeps the other tasks on
that CPU away from it.

% ------------------------------------------------------------------------------------------
\solution[fillaccent4]{c12:ex:5}{Priorities}{Design}

\xpart{a} Any two of these:
\begin{itemize}
  \item \textbf{Hardware interrupts.} A hard handler runs on top of whatever task is on its CPU,
    whatever its priority, and on \code{PREEMPT} so do the softirqs processed on the way out of it.
  \item \textbf{Anything of higher priority}: \code{SCHED_FIFO} and \code{SCHED_RR} tasks at 81 to
    99, every \code{SCHED_DEADLINE} task, since the deadline class ranks above the real-time one,
    and the kernel's stop-class threads, \code{migration/N}.
  \item \textbf{The throttle}, once real-time tasks on the CPU have used their 950 ms of the second
    (part c).
\end{itemize}
Note what does \emph{not}: on \code{PREEMPT_RT} interrupt threads run at \code{SCHED_FIFO} 50
(\file{kernel/sched/syscalls.c}:858), below 80. That includes the thread that services the device
the program is measuring, which is harmless only because the program blocks in \code{read}.

\xpart{b} Nothing of lower priority runs on that CPU again. A \code{SCHED_FIFO} task that never
blocks is never preempted by anything below it, so every ordinary task, including the shell you
would use to kill it, is starved, and the machine is hung for everything but interrupts and
higher-priority tasks. \textbf{What saves you is RT throttling.} By default real-time tasks may use
950,000 us of every 1,000,000 (\file{kernel/sched/rt.c}:15 and~21), so every second the CPU has
50 ms for everything else. The shell runs, slowly, and you can kill the task. 6.12 also has a
second safeguard with the same effect, the fair server, which reserves 50 ms in every 1,000 ms for
ordinary tasks (\file{kernel/sched/deadline.c}:1637).

\xpart{c} \textbf{In every \code{sched_rt_period_us}, one second by default, the real-time tasks on
a CPU may run for at most 950,000 us between them, so 50 ms of every second is left for everything
else.} The symptom: a real-time task running flat out stops dead for up to 50 ms once a second, a
periodic gap that nothing in the task explains, and the kernel prints ``\code{sched: RT throttling
activated}'' once (\file{kernel/sched/rt.c}:900).

\xpart{d} \textbf{When:} the real-time tasks' CPU use is bounded by design and shown to be,
because each of them blocks within its period by analysis and by measurement, and a 50 ms stall
once a second would itself break a deadline. A control loop on a core given over to it is the
usual case. It is done by writing $-1$ to \code{sched_rt_runtime_us}, which means ``no limit''
(\file{Documentation/scheduler/sched-rt-group.rst}:104). \textbf{What you owe the next person:} a
written record of where the setting is made and why; what bounds each real-time task's CPU use,
and the evidence; what now protects the system from a real-time task that loops, a hardware
watchdog or a higher-priority monitor; how to recover a board that hangs; and a test that fails if
the setting is lost or a real-time task spins. The safeguard is gone, and whoever changes the
real-time code next has to know that.

\xpart{e} It performed an \textbf{admission test}. It added the new task's runtime divided by its
period to that of every deadline task already admitted and compared the total with what the system
allows deadline tasks. By default that is the same 95\% of each CPU, since the \code{-rt} knobs are
used for deadline admission control
(\file{Documentation/scheduler/sched-deadline.rst}:594). When the total would be exceeded,
\code{sched_setattr} fails with \code{EBUSY} (\file{kernel/sched/syscalls.c}:687--689).
\code{SCHED_FIFO} has no admission control at all: it accepts any number of tasks at any priorities
and finds out at run time that they do not fit. \textbf{A refusal is more useful} because it happens
once, at configuration time, deterministically, where the program can report it and a test can catch
it. A missed deadline happens later, intermittently, under a load nobody reproduced, in the field.
On more than one CPU the test bounds how late a task can be rather than promising it never is
(\file{Documentation/scheduler/sched-deadline.rst}:453--463), and that is still a statement a
priority cannot make.

% ------------------------------------------------------------------------------------------
\solution[fillaccent4]{c12:ex:8}{What cyclictest is for}{Design}

\xpart{a} \code{cyclictest} runs one or more threads, normally at a real-time priority. Each sleeps
until an absolute time, notes when it actually woke, and sleeps again one interval later, and the
program reports the minimum, latest, average and maximum lateness of those wakeups for each thread.
It therefore measures how late the timer interrupt and the scheduler deliver a waiting task to its
CPU, under whatever load runs alongside it.

\xpart{b} The 45 us describes a timer's interrupt and the wakeup of a high-priority thread. Your
data travels a different path, none of which \code{cyclictest} ran: your device's line and the
interrupt controller, your handler, the step to its thread, your driver's locks, the wakeup of your
reader, and your reader's own priority. \code{cyclictest} usually runs at \code{SCHED_FIFO} 80 or
above, while a reader left at \code{SCHED_OTHER} waits for the fair scheduler, and on
\code{PREEMPT_RT} your interrupt thread at 50 waits behind any busy real-time task above it. On
\code{PREEMPT_RT} the timer that wakes a real-time \code{cyclictest} thread even expires in hard
interrupt context (\file{kernel/time/hrtimer.c}:2043--2044), so its path has no thread step at all.
Any stage of yours can be late while the timer path is fast.

\xpart{c} It covers the nearest thing to (1) and (3), \emph{for a timer}: the timer interrupt being
taken, and the woken thread getting the CPU, with no (2) for a real-time thread, as in b). It covers
none of the path for your device: not (1) for your line, not (2) from your handler to its thread,
not (3) from your driver's wakeup to your reader, and so not the sum.

\xpart{d} Any two of these:
\begin{itemize}
  \item \textbf{A GPIO and an instrument.} The handler, or the program in userspace, toggles a
    GPIO. An oscilloscope or logic analyser watches it against the signal that caused the
    interrupt, a signal generator or the device's own line, and the delay between the two edges is
    measured by something outside the system (\secref{c12:app:b:what-you-would-need-instead}).
  \item \textbf{A loopback on one clock.} The CPU drives a GPIO output wired to an interrupt input,
    reading its own counter as it does, and the handler reads the same counter. Both ends are then
    in one clock domain, which is what \code{qa-dev}'s register gives the course.
  \item \textbf{The kernel's tracers.} \code{irq_handler_entry} events and the \code{irqsoff},
    \code{preemptoff} and \code{wakeup_rt} tracers show where the time between handler and task
    goes. They cannot see the delay before the CPU takes the interrupt, which only an external
    reference can.
\end{itemize}

\xpart{e} Because the course runs on an emulator, and an emulated board has no pin to put a probe
on, so the techniques in d) are not available. A device that stamps the moment it raised its line
is the only reference point such a target can have, and it lets the technique, subtracting a
device's timestamp from the handler's, be taught to every reader on the same machine and checked
automatically by \code{make test}. A real device would bring whatever it happens to do and no
timestamp at all, since real hardware almost never has one (\secref{c12:app:b2}). The price is the
caveat in \secref{c12:app:b3}: the device's clock has its own origin, so the interval it gives is
only a spread.

% ------------------------------------------------------------------------------------------
\solution[fillaccent]{c12:ex:9}{The trade, measured}{Cross-check}

\xpart{a} The prediction is yours. The usual one is that \code{PREEMPT_RT} lowers both, and half of
it is wrong: what the course measured is in b).

\xpart{b} What the course measured and published in \secref{c12:app:a8}, one run of 2,000 samples
each, handler-to-userspace:

\begin{narrowtable}{@{}llrr@{}}
  \thead{Kernel} & \thead{Load} & \thead{Mean} & \thead{Worst} \\
  \midrule
  \code{PREEMPT} & idle & 449 us & 11,321 us \\
  \code{PREEMPT_RT} & idle & 551 us & 4,964 us \\
  \code{PREEMPT} & busy & 1,685 us & 29,729 us \\
  \code{PREEMPT_RT} & busy & 2,393 us & 18,559 us \\
\end{narrowtable}

\noindent Yours will differ, and \secref{c12:app:b6} says why that is expected.

\xpart{c}
\begin{itemize}
  \item \textbf{Which prediction:} with the table above, ``lower maximum'' is right both idle and
    busy, and ``lower mean'' is wrong both idle and busy.
  \item \textbf{The mean is worse} because \code{PREEMPT_RT} makes every path do more. Every
    interrupt costs a switch into its thread and out again, and softirqs run in threads too. Every
    \code{spinlock_t} acquisition is an \code{rt_mutex} operation that also does
    \code{rcu_read_lock} and \code{migrate_disable} (\file{kernel/locking/spinlock_rt.c}:50--51),
    and when contended it sleeps and does inheritance bookkeeping instead of spinning briefly. All
    of that is paid on every sample. \textbf{The maximum is better} because the long stretches
    during which nothing could preempt, \code{spinlock_t} sections anywhere in the kernel and softirq
    processing, are now preemptible, and inversion is bounded by inheritance. The reader's worst
    wait is no longer the longest non-preemptible section in the kernel. (In this lab, a handler
    requested with \code{devm_request_irq} is itself forced into a thread on \code{PREEMPT_RT}, so
    its first statement, the \code{ktime_get_ns()}, runs after the hard-interrupt-to-thread step;
    part of the thread cost the mean pays is outside the measured interval, which makes the rise in
    the mean more notable, not less.)
  \item \textbf{A colleague who measured only the mean} would conclude that \code{PREEMPT_RT} is
    worse, by 23\% idle ($551 / 449 = 1.23$) and 42\% loaded ($2{,}393 / 1{,}685 = 1.42$), and reject
    it. For a throughput system they would be right. For anything with a deadline they would be
    wrong, because a deadline is a statement about the maximum, which in the same runs fell by a
    factor of 2.3 idle ($11{,}321 / 4{,}964 = 2.28$) and 1.6 loaded ($29{,}729 / 18{,}559 = 1.60$).
  \item \textbf{In one sentence:} real time is defined by the worst case, and \code{PREEMPT_RT} buys
    a lower worst case by making the kernel preemptible almost everywhere, which costs something on
    every path, so a worse average with a better maximum is exactly what was paid for.
\end{itemize}

\xpart{d} The course published one repeat, of the idle \code{PREEMPT} run (\secref{c12:app:b5}):
11,321 us and then 2,502 us, a spread of 8,819 us, or a factor of $11{,}321 / 2{,}502 = 4.5$, on one
kernel with nothing changed. The difference between the kernels' idle maxima was
$11{,}321 - 4{,}964 = 6{,}357$ us, a factor of 2.3.
\begin{itemize}
  \item \textbf{No.} The run-to-run spread, 8,819 us, is larger than the kernel-to-kernel difference,
    6,357 us, and the repeat, 2,502 us, is below the \code{PREEMPT_RT} maximum of 4,964 us. One pair
    of runs could have been read either way.
  \item \textbf{Then you have established nothing about which kernel has the lower maximum.} What you
    have established is that the maximum of 2,000 samples is not a stable statistic on this target,
    which is itself the result \secref{c12:app:b5} states. The means point the same way idle and
    busy, but one run each supports no more than that.
  \item \textbf{How many:} enough that the two kernels' per-run maxima stop overlapping. If both came
    from the same distribution, the chance that all $n$ of one kernel's runs fall below all $n$ of
    the other's is $n!\,n!/(2n)! = 1/\binom{2n}{n}$. That is $1/20 = 5\%$ for three runs each,
    $1/70 \approx 1.4\%$ for four, and $1/252 \approx 0.4\%$ for five. Four or five runs of each, with
    no overlap, is what I would defend for the \emph{direction}. For a \emph{bound} no number of
    two-second runs is enough: that needs hours per run, under the product's load
    (\secref{c12:app:b5}).
\end{itemize}

\xpart{e} Everything in these numbers passed through far more than the guest kernel. QEMU translates
every guest instruction, so every path costs many times what it would on silicon. The interrupt
controller, the generic timer and \code{qa-dev} itself are device models, and \code{qa-dev}'s sample
timer is a QEMU timer that fires when QEMU gets round to it. QEMU's vCPUs are threads of an ordinary
process on the host, scheduled by a host kernel that is not real time and shares its machine with
everything else running there, so a vCPU descheduled on the host is a stall the guest can neither
see nor tell apart from its own latency; \secref{c12:app:b6} names that as frequently the largest
term. Inside the guest, both kernels carry \code{CONFIG_PROVE_LOCKING} and
\code{CONFIG_DEBUG_ATOMIC_SLEEP} (\file{kernel/qa.config}, \file{kernel/rt.config}), so every lock
acquisition does the validator's bookkeeping. \code{qa_latency_test} sets no scheduling policy, so
the reader is an ordinary \code{SCHED_OTHER} task, and under load it competes as an equal with four
processes that never block. The loaded figures are to a large extent the fair scheduler's time
slices, not interrupt latency. The measurement is also blind to parts of the path by construction.
Its absolute figure starts at the handler's \code{ktime_get_ns()}, so the time from the device
raising its line to the handler running appears only as a spread above an unknown minimum. On
\code{PREEMPT_RT} a handler requested with neither \code{IRQF_NO_THREAD} nor \code{IRQF_ONESHOT} is
forced into a thread (\file{kernel/irq/manage.c}:1368--1404), so its first statement runs in the
thread, and the step from hard interrupt to thread moves out of the absolute interval into the
spread: the two kernels' columns do not cover quite the same stretch of path. The device rewrites
its timestamp at every sample (\file{tools/qa-dev/qa-dev.c}:171), so a handler that runs more than
one period late reads a later sample's stamp, and the spread cannot show a delay longer than about
one period, which is exactly the kind that matters. And 2,000 samples are two seconds, too short to
have met the rare events a worst case is made of; one repeat moved the maximum by a factor of 4.5.
\textbf{So the absolute values cannot transfer.} Their dominant terms belong to the host and the
emulator, and none of those exists on the product; on real silicon the same kernel's figure would
typically be tens of microseconds (\secref{c12:app:b6}). \textbf{What survives} is what is common to
both columns and therefore cancels: the same emulator, host, load, reader and method on both
kernels, so the direction of the trade, a mean that rises under \code{PREEMPT_RT} and a maximum
that falls, is a property of the mechanism, and it held idle and busy alike. \textbf{What does not
survive} is every number, the factors of 2.3 and 1.6, the percentages, and even the direction of
the idle maximum until repeated runs separate the two kernels, since one repeat of a single kernel
moved further than the gap between them.

\xpart{f} The entry, for the kernel a deadline would choose:

\begin{aside}
\textbf{Handler-to-userspace latency}, from the \code{qa_latency} interrupt handler to a blocking
\code{read()} returning in a \code{SCHED_OTHER} process: mean 551 us, maximum 4,964 us, idle; mean
2,393 us, maximum 18,559 us, under four CPU-bound processes. Linux 6.12.30 \code{PREEMPT_RT}, arm64
defconfig with \file{kernel/qa.config} and \file{kernel/rt.config}, including
\code{CONFIG_PROVE_LOCKING} and \code{CONFIG_DEBUG_ATOMIC_SLEEP}. \code{qemu-system-aarch64} 9.2.0,
\code{-M virt -cpu cortex-a72 -smp 2 -m 512}, software emulation, on a named host with its load
stated. Device period 1 ms; 2,000 samples per condition, about two seconds; one run each, not
repeated.
\end{aside}

\noindent The sentence: \emph{``This is not interrupt-to-userspace latency: it omits the time from
the device asserting its interrupt to the handler running, which this target can report only as a
spread above an unknown minimum; and because it was measured under an emulator whose host
scheduling dominates it, it predicts nothing about this kernel on real hardware.''} Written out
with every condition, the entry shows that a figure from this target belongs in a lab report and
never in a product datasheet. For a datasheet you measure the product, as in
\secref{c12:app:b:what-you-would-need-instead}.

\xpart{g} One defensible answer: \textbf{a latency measurement on the product's own hardware}, with
debug options off, for hours, under the product's load, against an external reference. Its latency
has been characterised only inside an emulator, and \secref{c12:app:b6} is where this book says why
that characterisation predicts nothing about silicon. The first place the course made the same
point about a number and its conditions was Chapter~\ref{c2:ch}'s Cross-check,
Exercise~\ref{c2:ex:9}. An answer that names concurrency testing instead is as good if it argues
why that comes first. Every test in the course runs one thing at a time, and lockdep finds only the
orderings your tests reach (\secref{c7:app:b5}), so a driver whose locking has never been exercised
under contention is unverified in a way no latency figure touches.
